%% kapa: exam=1 round=29 year=2018 period=1 form=A subject=civil_law subject_ko=민법 seq=18 number=18
%% source: Q-Net 감정평가사 기출문제 (cst00309) artlSeq=5202808 file=extracted/1008/5202808/2018년 제29회 감정평가사 1차시험 기출문제/2018년 제29회 감정평가사 1차시험 기출문제/1교시 A형.pdf
%% conversion: kordoc 4.24.51 -> tools/gen1.py (auto); review=proofread-sonnet; flags=none
\begin{question}{18}
법률행위의 조건에 관한 설명으로 옳은 것은? (다툼이 있으면 판례에 따름)
\begin{choices}
  \item 조건의 성취가 미정인 권리는 일반규정에 의하여 처분, 상속할 수 있으나 담보로 제공할 수는 없다.
  \item 조건이 법률행위의 당시 이미 성취한 것인 경우에는 그 조건이 해제조건이면 조건없는 법률행위로 한다.
  \item 조건의 성취로 인하여 이익을 받을 당사자가 신의성실에 반하여 조건을 성취시킨 때에도 상대방은 그 조건이 성취하지 아니한 것으로 주장할 수 없다.
  \item 조건부 법률행위에 있어 조건의 내용 자체가 불법적인 것이어서 무효일 경우 그 조건만을 분리하여 무효로 할 수 있다.
  \item 조건의 성취로 인하여 불이익을 받을 당사자가 신의성실에 반하여 조건의 성취를 방해한 경우, 조건이 성취된 것으로 의제되는 시점은 신의성실에 반하는 행위가 없었더라면 조건이 성취되었으리라고 추산되는 시점이다.
\end{choices}
\end{question}
\answer{5}
